% Abstract — no \section command, no citations, no math, no special characters.
% IEEE hard rule: keep under 150 words for conference papers.

% Kenya's livestock industry accounts for over 30\% of all agricultural products and contributes 42\% of the country's agricultural GDP and 12\% of its GDP. 
% It is thus essential to Kenya's economy as it promotes economic growth, poverty reduction, job creation, sustainable livelihoods, and food security. Quantitative livestock sector planning in Sub-Saharan Africa is frequently constrained by the absence of reproducible data driven analytical models that bridges biological parameters and environmental factors. This study covers a forecasting framework that includes nine of Kenya's principal livestock value chains; dairy cattle, beef cattle, camels, goats, sheep, poultry, pigs, rabbits, and apiculture initialized using 2019 as the base year, with projections simulated quarterly from 2020 to 2041. Drawing on the 2019 KNBS livestock census as the baseline, the model applies a quarterly transition sequence, covering mortality, culling, offtake, age-based maturation, and reproduction, across all cohorts for each value chain (including calves, heifers, kids, weaners, boars etc.). The parameters used are sourced from Kenya-specific and Sub-Saharan African peer-reviewed literature. Production modules were linked to each population model translating to herd and flock dynamics into species-specific output metrics: milk for dairy cattle and camels, meat and carcass weight for beef cattle and small ruminants, eggs and broiler meat for poultry, and honey for apiculture. A carrying capacity constraint was incorporated to represent environmental limits on herd growth, based on forage availability estimated from satellite-derived vegetation indices to capture seasonal productivity, and adjusted by a utilization factor. Herd size was then constrained through a dynamic balance between feed supply and herd intake requirements. Lastly, the model links production outputs to a feed demand module estimating dry matter requirements by species and agro-ecological zone. Outputs from the model show a positive CAGR growth ranging from 3.6\% to 6\% across all value chains demonstrating the robustness of a biologically grounded modelling approach, where coherent parameterisation of herd dynamics, productivity factors, and system constraints enables the generation of defensible forecasts. The model not only captures growth trajectories but also reveals underlying productivity bottlenecks and structural inefficiencies within the sector. Beyond forecasting, the model offers a comprehensive analytical framework for assessing the relationships between livestock populations, output volumes, and resources required. This makes it possible to identify important pressure points that could go unnoticed in aggregate assessments, such as feed deficiencies, low off-take rates, and subpar reproductive performance. The insights from the model can guide strategic interventions in food security planning, direct targeted investments in infrastructure and influence data-driven livestock policies.

% shortened the abstract to under 200 words below:

Kenya's livestock sector contributes 42\% of agricultural GDP and 12\% of national GDP, yet quantitative planning remains constrained by the absence of reproducible, data-driven analytical models. This study presents a cohort-based quarterly forecasting framework spanning nine value chains: dairy cattle, beef cattle, camels, goats, sheep, poultry, pigs, rabbits, and apiculture, initialized from the 2019 KNBS livestock census with projections from 2020 to 2041. The model applies a quarterly transition sequence covering mortality, culling, offtake, age-based maturation, and reproduction across all cohorts. Production modules translate herd and flock dynamics into species-specific output metrics, while a carrying capacity constraint — derived from satellite-based vegetation indices and adjusted by a utilization factor — regulates growth against forage availability. A feed demand module estimates dry matter requirements by species and agro-ecological zone. Projected outputs indicate a CAGR of 3.6-6\% across all value chains. Beyond forecasting, the framework exposes productivity bottlenecks and structural inefficiencies, including feed deficits, low offtake rates, and suboptimal reproductive performance, that aggregate assessments typically obscure. The model provides an analytical basis for food security planning, targeted infrastructure investment, and evidence-based livestock policy.